% SPDX-License-Identifier: CC-BY-SA-4.0
% Author: Matthieu Perrin

\begingroup

\SetKwFunction{Propose}{propose}
\SetKwFunction{Filter}{filter}
\SetKwBroadcast{MB}{mb}
\SetKwMessage{FILTER}{filter}
\SetKwData{F}{f}
\SetKwVariable{Received}{received}
\newcommand\Adopt{\textup{\textsc{adopt}}}
\newcommand\Abort{\textup{\textsc{abort}}}
\newcommand\Commit{\textup{\textsc{commit}}}

\begin{exercice}[Abort-Adopt-Commit]

  On considère l'abstraction \emph{Abort-Adopt-Commit} (AAC).
  Chaque processus propose une valeur de $\{\False,\True\}$ et retourne une valeur parmi :
  $\Commit(\False), \Adopt(\False), \Abort, \Adopt(\True), \Commit(\True)$.

  \begin{definition}[Abort-Adopt-Commit]
    Une instance d'Abort-Adopt-Commit vérifie les propriétés suivantes :
    \begin{description}
    \item[AAC-Validité.]
      Si tous les processus proposent la même valeur $v$, alors
      seule $\Commit(v)$ peut être retournée.
    \item[AAC-Quasi-Accord.]
      Les valeurs retournées par deux processus sont identiques ou adjacentes
      dans la liste $\Commit(\False), \Adopt(\False), \Abort, \Adopt(\True), \Commit(\True)$.
    \item[AAC-Terminaison.]
      Toute invocation de $\Propose$ par un processus correct termine.
    \end{description}
  \end{definition}
  
  On considère maintenant l'algorithme~\ref{algo:Filter}, qui implémente l'abstraction $\Filter$.

  \begin{algorithm}[H]
    \lLVariables{\At $p_i$}{
      $\Received_i \leftarrow \emptyset$\tcp*[f]{Ensemble des valeurs $v$ pour lesquelles $p_i$ a reçu $\FILTER(v)$}
    }
    \Operation{$\Filter(v)$ \At $p_i$}{
      \nl \MB.\Broadcast $\FILTER(v)$\;
      \nl \Return $\Received_i$
    }
    \nl \lUpon{\MB.\Deliver $\FILTER(v)$ \From $p_j$ \At $p_i$}{
      $\Received_i \leftarrow \Received_i \cup \{v\}$
    }
    \caption{Abstraction \textsc{filter}}
    \label{algo:Filter}
  \end{algorithm}

  \begin{questions}

  \item Montrer que si tous les processus proposent la même valeur $v$,
    alors $\{v\}$ est la seule valeur qui peut être retournée par $\Filter$.

  \item Quelles valeurs une invocation de $\Filter$ peut-elle retourner si les valeurs proposées sont booléennes ?

  \item Caractériser les valeurs que peuvent retourner deux processus dans une même exécution.
    En particulier :
    \begin{enumerate}
    \item donner une exécution dans laquelle un processus retourne $\{\False\}$
      et un autre $\{\False, \True\}$ ;
    \item montrer qu'il est impossible que deux processus retournent
      respectivement $\{\False\}$ et $\{\True\}$.
    \end{enumerate}

  \item On considère maintenant deux instances partagées de \Filter, $\F_1$ et $\F_2$.
    Chaque processus $p_i$ appelle d'abord $\F_1.\Filter(v_i)$, puis appelle
    $\F_2.\Filter(S_i)$ avec l'ensemble $S_i$ retourné par $\F_1$.
    Caractériser les valeurs que peut retourner $\F_2$, ainsi que les couples
    de valeurs que peuvent retourner deux processus dans une même exécution
    de $\F_2.\Filter(\F_1.\Filter(v))$.

  \item En déduire une implémentation d'Abort-Adopt-Commit à partir des deux
    instances $\F_1$ et $\F_2$ de l'abstraction $\Filter$.
    Préciser comment chaque valeur retournée par $\F_2$ est interprétée,
    puis justifier les propriétés \textsc{AAC-Validité},
    \textsc{AAC-Quasi-Accord} et \textsc{AAC-Terminaison}.
    
  \end{questions}

\end{exercice}

    

%  \item On souhaite utiliser l'abstraction $\Filter$ pour implémenter
%    Abort-Adopt-Commit.
%    Expliquer pourquoi la transformation
%    \(      \{v\} \longmapsto \textsc{commit}(v)
%    \qquad\text{et}\qquad
%    \{0,1\} \longmapsto \textsc{abort}
%    \)
%    ne vérifie pas la propriété \textsc{AAC-Quasi-Accord}.

%\begin{exercice}[\footnote{[B83] Michael Ben-Or. \emph{Another advantage of free choice for completely asynchronous agreement protocols.} PODC 1983}]
%  
%  Le but de cet exercice est d'étudier l'algorithme de consensus randomizé de Ben-Or. Il s'agit d'un algorithme de consensus binaire,
%  c'est-à-dire tel que les seules valeurs qui peuvent être proposées sont $0$ et $1$.
% 
%  \begin{definition}[Abort-Adopt-Commit]
%    Abort-Adopt-Commit est un affaiblissement du consensus dans lequel tous les processus peuvent
%    proposer $0$ ou $1$, et retourner $\textsc{abort}$, $\textsc{adopt}(v)$ ou $\textsc{commit}(v)$, où $v \in \{0, 1\}$, tel que :
%    \begin{description}
%    \item [AAC-Validité.] Si tous les processus proposent la même valeur $v$, $\textsc{commit}(v)$ est la seule sortie possible.
%    \item [AAC-Quasi-Accord.] Si deux processus retournent des valeurs différentes, alors ces valeurs sont adjacentes dans la liste : 
%      $\textsc{commit}(0); \textsc{adopt}(0); \textsc{abort}; \textsc{adopt}(1); \textsc{commit}(1)$.
%    \item [AAC-Terminaison.] Une invocation de ${\sf propose}()$ par un processus correct termine.
%    \end{description}
%  \end{definition}
%  
%  Dans un premier temps, on s'intéresse à l'algorithme~\ref{algo:Filter} ci-dessous.
%  
%  \begin{algorithm}
%    \SetKwFunction{Filter}{filter}
%    \SetKwFunction{Local}{local\_coin}
%    \SetKwFunction{Rand}{rand}
%    \lLVariables{}{$\mathit{received}_i \leftarrow \emptyset$}
%    \Operation{$\Filter(v)$}{
%      \SBroadcast{mutual} $\textsc{filter}(v)$\;
%      \Return $\mathit{received}_i$\;
%    }
%    \lWhen{\Deliver{mutual} $\textsc{filter}(v)$ \From $p_j$}{
%      $\mathit{received}_i \leftarrow \mathit{received}_i \cup \{v\}$%
%    }
%    \caption{Algorithme filter dans le modèle $CAMP_{n,t}[mutual-broadcast]$}
%    \label{algo:Filter}
%  \end{algorithm}
% 
%  \begin{questions}
% 
%  \item Lorsque les valeurs proposées appartiennent à ${0,1}$, caractériser les
%    ensembles qui peuvent être retournés par $\Filter$.
%    Caractériser également les couples de résultats qui peuvent être retournés
%    par deux processus dans une même exécution.
% 
%  \item Que peut-on conclure si tous les processus proposent la même valeur $v$ ?
% 
%  \item On souhaiterait transformer les résultats de $\Filter$ en valeurs
%    $\textsc{abort}$, $\textsc{adopt}(v)$ et $\textsc{commit}(v)$.
%    Expliquer pourquoi associer directement
%    ${v}$ à $\textsc{commit}(v)$ et ${0,1}$ à $\textsc{abort}$
%    ne satisfait pas la propriété d'\textsc{AAC-Quasi-Accord}.
% 
%  \item Proposer, à partir de l'abstraction $\Filter$, une implémentation
%    d'Abort-Adopt-Commit et justifier qu'elle en vérifie les trois propriétés.
% 
%  \end{questions}
% 
%\end{exercice}

\endgroup
\endinput
